\chapter{The identity component}
\label{mk-ch-components}

\begin{definition}[Degree components]
  \label{mk-def-picard-degree}
  \dcref{2.2,custom}
  \uses{mk-thm-picard-representability}
  The component \(\operatorname{Pic}^d_{C/k}\) is the locus of line bundles of
  degree \(d\).  Tensor product induces
  \[
    \operatorname{Pic}^d_{C/k}\times\operatorname{Pic}^e_{C/k}
    \longrightarrow\operatorname{Pic}^{d+e}_{C/k}.
  \]
\end{definition}

\begin{theorem}[Identity component]
  \label{mk-thm-identity-component}
  \dcref{5.1}
  \uses{mk-thm-picard-representability,mk-prop-picard-group}
  \cite[Lem.~5.1]{kleiman-picard}
  If \(G\) is a group scheme locally of finite type over a field, its connected
  component \(G^0\) containing the identity is an open and closed subgroup
  scheme, geometrically irreducible, and formation of \(G^0\) commutes with
  extension of the ground field.
\end{theorem}

\begin{theorem}[The Picard variety]
  \label{mk-thm-pic0-abelian}
  \dcref{5.4,5.11,5.13,5.14,custom}
  \uses{mk-thm-identity-component,mk-prop-genus-tangent,mk-def-genus}
  \cite[Thms.~5.4 and 5.11, Cors.~5.13--5.14]{kleiman-picard}
  The identity component
  \[
    \operatorname{Pic}^0_{C/k}:=(\operatorname{Pic}_{C/k})^0
  \]
  is a smooth, proper, geometrically connected commutative group scheme over
  \(k\), hence an abelian variety, and
  \[
    \dim\operatorname{Pic}^0_{C/k}=g(C).
  \]
\end{theorem}
